\subsection{BASES OF A VECTOR SPACE}

THEOREM 1. Every vector space over a field K is a free K-module.

It must be proved that every vector space admits a basis; this will follow from the following more precise theorem:

THEOREM 2. Given a generating system S of a vector space E over a field K and a free subset L of E contained in S, there exists a basis B of E such that \(L \subset B \subset S\) .

Theorem 1 follows from this statement by taking \(\mathbf{L} = \varnothing\) .

To prove Theorem 2, we note that the set \(\mathfrak{L}\) of free subsets of E contained in S, ordered by inclusion, is an inductive set (Set Theory, III, §2, no. 4), by virtue of §1, no. 11; so is the set \(\mathfrak{M}\) of free subsets containing L and contained in S. By Zorn's Lemma, \(\mathfrak{M}\) admits a maximal element B and it suffices to prove that the vector subspace of E generated by B is equal to E. This follows immediately from the definition of B and the following lemma:

Lemma 1. Let \((a_{i})_{i\in I}\) be a free family of elements of \(\mathbf{E}\) ; if \(b\in \mathbf{E}\) does not belong to the subspace \(\mathbf{F}\) generated by \((a_{i})\) , the subset of \(\mathbf{E}\) consisting of the \(a_{i}\) and \(b\) is free.

Suppose that there were a relation \(\mu b + \sum_{i} \lambda_i a_i = 0\) with \(\mu \in K\) and \(\lambda_i \in K\) for all \(i \in I\) , the family \((\lambda_i)\) having finite support; if \(\mu \neq 0\) , it would follow that \(b = -\sum_{i} (\mu^{-1} \lambda_i) a_i\) and hence \(b \in F\) contrary to the hypothesis; hence \(\mu = 0\) and the relation becomes \(\sum_{i} \lambda_i a_i = 0\) , which implies \(\lambda_i = 0\) for all \(i \in I\) by hypothesis; whence the lemma.

COROLLARY. For a subset B of a vector space E, the following properties are equivalent:

\begin{enumerate}
\def\labelenumi{(\alph{enumi})}
\item
  B is a basis of E.
\item
  B is a maximal free subset of E.
\item
  B is a minimal generating system of E.
\end{enumerate}

This follows immediately from Theorem 2.

Example. Given a ring A and a subfield K of A, A is a (right or left) vector space over K and therefore admits a basis; in particular, every extension field of a field K has a basis as a left (resp. right) vector space over K. *Thus the field R of real numbers admits an (infinite) basis as a vector space over the field Q of rational numbers; such a basis of R is called a Hamel basis.*

Remark. For a family \((a_{i})_{i\in I}\) of elements of a vector space E over a field K to be free, it is necessary and sufficient that, for all \(k\in I\) , \(a_{k}\) belong to no subspace of E generated by the \(a_{i}\) of index \(i\neq k\) . We know that this condition is necessary in any module (§1, no. 11, Remark 1). It is sufficient by virtue of Lemma 1, as is seen immediately arguing by reductio ad absurdum and considering a minimal related subfamily of \((a_{i})\) .

